\section{Problem Setting}
\label{sec:problem-setting}

The analysis distinguishes three objects that coincide only in idealized
settings: the causal feature of an infinite input history, the implementable
feature obtained from a finite window, and the empirical sample of dependent
windows used to train the readout. This section fixes these objects before the
quantum architecture is specified.

\subsection{Temporal data and windowed observations}

Let \(Z=(Z_t)_{t\in\Z}\), with \(Z_t=(X_t,Y_t)\), be a stationary stochastic
process on \(\R^d\times\R\). We write
\(x_{-\infty:0}=(\ldots,x_{-1},x_0)\) for a semi-infinite input history and
\[
    x_{-w+1:0}=(x_{-w+1},\ldots,x_0)\in(\R^d)^w
\]
for its most recent \(w\) inputs. A temporal predictor maps a history to a
scalar. The target need not have finite memory; a finite window is instead an
implementable approximation whose error will later be controlled by reservoir
contraction.

For a window length \(w\geq1\) and stride \(s=w+g\), with \(g\geq0\), define
\begin{equation}
    W_j
    :=
    \bigl(X_{js-w+1:js},Y_{js}\bigr),
    \qquad j\in\Z.
    \label{eq:window-process}
\end{equation}
The observed sample is \(\cD_N=(W_1,\ldots,W_N)\). The gap \(g\) controls the
raw-time separation between consecutive windows, but nonoverlap alone does not
make them independent.

\subsection{Causal and finite-window reservoir features}

There are \(R\) independently parameterized sub-reservoirs. Sub-reservoir
\(r\) has an \(n\)-qubit state \(\rho_t^{(r)}\) and an input-dependent channel
\(T_{r,x}\). For a semi-infinite history, its causal state at time zero is
written \(\rho_{\infty,0}^{(r)}(x_{-\infty:0})\). For a finite window, the
recursion starts at time \(-w\) from a fixed state
\(\rho_{\mathrm{init}}^{(r)}\), producing
\(\rho_{w,0}^{(r)}(x_{-w+1:0})\) after the last input is consumed.

Let \(\cO\) be a finite family of Hermitian observables with
\(\opNorm{O}\leq1\). The exact infinite-history and finite-window feature maps
are
\begin{align}
    \PhiInf(x_{-\infty:0})
    &:=
    \bigl(
      \Tr[O\rho_{\infty,0}^{(r)}]
    \bigr)_{r\in[R],\,O\in\cO},
    \label{eq:infinite-feature}\\
    \PhiWin(x_{-w+1:0})
    &:=
    \bigl(
      \Tr[O\rho_{w,0}^{(r)}]
    \bigr)_{r\in[R],\,O\in\cO}
    \in\R^{p},
    \qquad p=R|\cO|.
    \label{eq:window-feature}
\end{align}
These vectors concatenate expectations from a tuple of independently evolved
states; no tensor-product state is introduced for multiplexing. A finite-shot
estimate \(\widehat\Phi_w\) is a distinct random object and will be treated
separately from the exact feature map.

\subsection{Kernel readout and the three risks}

Let \(\kappa:\R^p\times\R^p\to\R\) be a normalized Mat\'ern kernel and
\(\cH_\kappa\) its RKHS. For a radius \(\Lambda>0\), define
\[
    \Hball
    :=
    \set{h\in\cH_\kappa:\norm{h}_{\cH_\kappa}\leq\Lambda}.
\]
Normalization gives \(\kappa(z,z)=1\), and hence
\begin{equation}
    \abs{h(z)}
    \leq\norm{h}_{\cH_\kappa}\sqrt{\kappa(z,z)}
    \leq\Lambda,
    \qquad h\in\Hball.
    \label{eq:prediction-bound}
\end{equation}

For squared loss \(\ell(\widehat y,y)=(\widehat y-y)^2\), the causal target is
the infinite-history risk
\begin{equation}
    \risk_\infty(h)
    :=
    \E\!\left[
      \ell\bigl(h(\PhiInf(X_{-\infty:0})),Y_0\bigr)
    \right].
    \label{eq:infinite-risk}
\end{equation}
The finite-window approximation is
\begin{equation}
    \risk_w(h)
    :=
    \E\!\left[
      \ell\bigl(h(\PhiWin(X_{-w+1:0})),Y_0\bigr)
    \right],
    \label{eq:window-risk}
\end{equation}
and the quantity optimized from the observed windows is
\begin{equation}
    \erisk_{N,w}(h)
    :=
    \frac1N\sum_{j=1}^N
    \ell\bigl(h(\PhiWin(X_{js-w+1:js})),Y_{js}\bigr).
    \label{eq:empirical-risk}
\end{equation}
The learning analysis will therefore have two parts: a statistical comparison
between \(\risk_w\) and \(\erisk_{N,w}\), and a deterministic comparison between
\(\risk_\infty\) and \(\risk_w\).

\subsection{Dependence and analytical scope}
\label{sec:analytical-scope}

For sub-sigma-algebras \(\mathcal A\) and \(\mathcal B\), let
\(\beta(\mathcal A,\mathcal B)\) denote their absolute-regularity coefficient.
We use
\begin{equation}
    \beta_Z(q)
    :=
    \sup_{t\in\Z}
    \beta\!\left(
      \sigma(Z_u:u\leq t),
      \sigma(Z_u:u\geq t+q)
    \right),
    \qquad q\geq0,
    \label{eq:beta-definition}
\end{equation}
and call the process beta mixing when \(\beta_Z(q)\to0\).

The primary learning result is first stated for a fixed design. The
projection, reservoir parameters, initialization states, observable family,
window scheme, and Mat\'ern parameters are chosen independently of the sample
on which that result is evaluated; when specialized to kernel ridge regression,
its regularization parameter is fixed on the same basis. Random design elements,
such as the JL matrix or reservoir seed, may be drawn independently and then
conditioned upon. We assume that \(x\mapsto V_r(x)\), equivalently
\(x\mapsto T_{r,x}\), is measurable, that the resulting loss classes are
measurable (or admit the usual separable versions), that
\(\abs{Y_t}\leq\Upsilon\) almost surely for some \(\Upsilon>0\), and that the
Mat\'ern kernel is normalized. Results using its spectral second moment require
\(\nu>1\). The main risk bound concerns exact observable expectations and is
uniform over \(\Hball\). \Cref{cor:finite-family-selection} later extends this
uniformity to sample-dependent selection from a finite family of configurations
that is fully declared before the sample is examined, including all random
seeds and preprocessing maps. Continuous or adaptively generated
hyperparameter families remain outside that extension.

With the data, feature maps, risks, and dependence structure in place, we next
specify the concrete \textsc{QuaRK} pipeline that produces \(\PhiWin\).
